\section{Transfer to policy cost}\label{sec:transfer-limits}
Preserving teacher-state cost comparisons does not certify an updated
policy. Summed advantages must be evaluated on student trajectories;
approximate continuation additionally requires sensitivity and curvature
control. The bounds in \ref{app:transfer} separate target gain, fitting
error, sensitivity discrepancy and continuation remainder. Local
improvement sets do not remove the shared-policy constraint: an exact
three-context example has nonempty sets with no common student,
although average-cost improvement is attainable (\ref{sec:sets}).

Separate nonlinear studies test these limits. In 50 continuation-audit
fits, the guided policy reduces cost relative to its initialization by
7.74\% and 18.64\% on mechanics and reaction tasks, but remains 0.42\%
and 2.60\% above its teacher. The comparison does not establish an
advantage over adaptive regularization. Sampled fitting costs can exceed local target gains,
and 26 of 4,800 perturbed-action checks violate the empirical curvature
bound (\ref{app:audit-protocol}--\ref{app:audit-results}).

A 56-fit improvement-set study retains an independent five-seed
comparison: mean cost is 0.54\% below point-target fitting, with a
paired interval crossing zero, and 1.28\% above upper-model fitting.
All five selected set policies remain worse than their teacher
(\ref{sec:set-experiments}). These different procedures neither
establish general optimization superiority nor identify the same
causal mechanism as the controlled loss intervention.
